\documentclass[12pt,a4paper]{article}

\usepackage{amsmath}
\usepackage{amsfonts}
\usepackage{amsthm}
\usepackage{mathtools}
\PassOptionsToPackage{hyphens}{url}\usepackage{hyperref}

% Personal Desktop
\newcommand{\pathcommon}[1] {F:/latex_commons/#1}

% Heriot-Watt PC
% \newcommand{\pathcommon}[1] {C:/Users/jc2021/Documents/GitHub/latex_commons/#1}

\input{\pathcommon{standard_theorems.tex}}
\input{\pathcommon{derivatives.tex}}

\newcommand{\dobibliography} {\bibliographystyle{apalike} \bibliography{\pathcommon{bibliography}}}

\newcommand{\realfunction}[2] {{#1: #2 \rightarrow \mathbb{R}}}
\newcommand{\cinfty} {\mathcal{C}^{\infty}}
\newcommand{\analytic} {\mathcal{C}^{\omega}}

\title{Integral Representation of Increasing Log-Convex-Shaped Functions via Exponential Softplus}
\author{Juan Casanova}

\begin{document}

\maketitle

This is the abstract.

\tableofcontents

Maybe an intro here about why this might be useful and so on.

\section{Background: Log-convexity}

The definitions and results in this section are classical and are introduced here to provide relevant context to readers.

\begin{definition}[Log-convex function]
A real function $\realfunction{f}{[a,b]}$ is log-convex if $\log(f)$ is convex. Formally, log-convexity can be characterized by any of the following two equivalent constraints, for all $x_1,x_2 \in [a,b]$ and all $t \in [0,1]$.
\begin{equation}
\log(f(t x_1+(1-t) x_2)) \leq t \log(f(x_1)) + (1-t) \log(f(x_2))
\end{equation}
\begin{equation}
f(t x_1 + (1-t) x_2) \leq f(x_1)^{t} f(x_2)^{1-t}
\end{equation}
\end{definition}

Typical examples of log-convex functions include exponential polynomials such as $e^x$ or $e^{x^2}$.\\

INCLUDE CHART OF EXPONENTIALS HERE

In twice differentiable functions, convexity and log-convexity can be characterized by constraints on the derivatives.

\begin{lemma}
Let $\realfunction{f}{[a,b]}$ be a twice differentiable function. Then, $f$ is convex if and only if $\secondderiv{f} \geq 0$ in all of $[a,b]$.
\end{lemma}

\begin{lemma}
Let $\realfunction{f}{[a,b]}$ be a twice differentiable function. Then, $f$ is log-convex if and only if $\secondderiv{f}f \geq (\firstderiv{f})^2$ in all of $[a,b]$.
\end{lemma}

\begin{proof}
$f$ is log-convex if and only if $\log f$ is convex if and only if $\secondderiv{(\log f)} \geq 0$.
$$\firstderiv{(\log f)} = \frac{\firstderiv{f}}{f}$$
$$\secondderiv{(\log f)} = \frac{\secondderiv{f}f - \firstderiv{f}\firstderiv{f}}{f^2} = \frac{\secondderiv{f}f - (\firstderiv{f})^2}{f^2}$$

But since $f^2 \geq 0$ in $[a,b]$, we have that $\secondderiv{(\log f)} \geq 0$ if and only if:
$$\secondderiv{f}f - (\firstderiv{f})^2 \geq 0$$ if and only if
$$\secondderiv{f}f \geq (\firstderiv{f})^2$$
\end{proof}

An important function to understand the structure of log-convex functions is the log-sum-exp function:

\begin{definition}{log-sum-exp}
\begin{equation}
LSE(x_1,...,x_n) = \log(e^{x_1}+...+e^{x_n})
\end{equation}
\end{definition}

\begin{lemma}[Convexity of log-sum-exp]
log-sum-exp is a convex function.
\end{lemma}

A proof and more details on the log-sum-exp function can be found in textbooks on convex functions, such as \cite{2004_boyd_convex_optimization}.\\

Log-convex functions form a convex cone. That is, they are closed under multiplication by nonnegative scalars and addition:

\begin{lemma}
Let $f, g$ be log-convex functions, and let $\alpha, \beta \geq 0$ be nonnegative scalars. Then $\alpha f + \beta g$ is log-convex.
\end{lemma}

\begin{proof}
Multiplication by nonnegative scalars can be shown arithmetically from the definition of log-convexity:
$$(\alpha f)(t x_1 + (1-t) x_2) = \alpha f(t x_1 + (1-t) x_2)$$
but since $f$ is log-convex it follows that
\begin{equation*}
\begin{array}{l}
\alpha f(t x_1 + (1-t) x_2) \leq \alpha f(x_1)^t f(x_2)^{1-t} = (\alpha^t \cdot \alpha^{1-t}) f(x_1)^t f(x_2)^{1-t} =\\
= (\alpha f(x_1))^t (\alpha f(x_2))^{1-t} = ((\alpha f)(x_1))^t ((\alpha f)(x_2))^{1-t}.
\end{array}
\end{equation*}

Addition can be shown from the convexity of log-sum-exp. Let $f,g$ be log-convex functions. Then $\log f$ and $\log g$ are convex. Consider $\log (f+g) = \log(e^{\log f} + e^{\log g}) = LSE(\log f, \log g)$. Composition of convex functions is convex, and therefore $\log (f+g)$ is convex, and thus $f+g$ is log-convex.
\end{proof}

\section{Log-convex-shaped functions}

Log-convexity is not a shape constraint, in the sense that it is not preserved by vertical translation.

\section{Integral representation of increasing convex functions}

Traditional ramp decomposition.

\section{Integral representation of increasing log-convex-shaped functions}

Push the decomposition properly to log-convex-shaped functions.

\dobibliography

\end{document}